% !TeX root = ../sections/03_solutions.tex
\subsection{1.1.B.2.1}
\begin{proof}
    まず，
    \[
        \lim_{n\to\infty}
        \sqrt{n}\left(\sqrt{n+1}-\sqrt{n}\right)
        =
        \frac{1}{2}
    \]
    を示す。

    左辺の数列を変形する。
    有理化すると，
    \[
        \sqrt{n}\left(\sqrt{n+1}-\sqrt{n}\right)
        =
        \sqrt{n}\cdot
        \frac{(n+1)-n}{\sqrt{n+1}+\sqrt{n}}
    \]
    である。したがって
    \[
        \sqrt{n}\left(\sqrt{n+1}-\sqrt{n}\right)
        =
        \frac{\sqrt{n}}{\sqrt{n+1}+\sqrt{n}}
    \]
    となる。

    分母分子を $\sqrt{n}$ で割ると，
    \[
        \frac{\sqrt{n}}{\sqrt{n+1}+\sqrt{n}}
        =
        \frac{1}{\sqrt{1+\frac{1}{n}}+1}
    \]
    である。

    よって，
    \[
        \sqrt{n}\left(\sqrt{n+1}-\sqrt{n}\right)
        =
        \frac{1}{\sqrt{1+\frac{1}{n}}+1}
    \]
    である。

    ここで，任意に $\epsilon>0$ を取る。
    示したいのは，十分大きな $n$ に対して
    \[
        \left|
            \frac{1}{\sqrt{1+\frac{1}{n}}+1}
            -
            \frac{1}{2}
        \right|
        <\epsilon
    \]
    が成り立つことである。

    実際，
    \[
        \left|
            \frac{1}{\sqrt{1+\frac{1}{n}}+1}
            -
            \frac{1}{2}
        \right|
        =
        \left|
            \frac{2-\left(\sqrt{1+\frac{1}{n}}+1\right)}
            {2\left(\sqrt{1+\frac{1}{n}}+1\right)}
        \right|
    \]
    であるから，
    \[
        =
        \frac{\sqrt{1+\frac{1}{n}}-1}
        {2\left(\sqrt{1+\frac{1}{n}}+1\right)}
    \]
    となる。

    さらに，
    \[
        \sqrt{1+\frac{1}{n}}-1
        =
        \frac{\frac{1}{n}}{\sqrt{1+\frac{1}{n}}+1}
    \]
    であるから，
    \[
        \left|
            \frac{1}{\sqrt{1+\frac{1}{n}}+1}
            -
            \frac{1}{2}
        \right|
        =
        \frac{1}
        {2n\left(\sqrt{1+\frac{1}{n}}+1\right)^2}
    \]
    を得る。

    ここで
    \[
        \sqrt{1+\frac{1}{n}}\geq 1
    \]
    であるから，
    \[
        \sqrt{1+\frac{1}{n}}+1\geq 2
    \]
    である。したがって
    \[
        \left(\sqrt{1+\frac{1}{n}}+1\right)^2\geq 4
    \]
    である。

    よって
    \[
        \left|
            \frac{1}{\sqrt{1+\frac{1}{n}}+1}
            -
            \frac{1}{2}
        \right|
        \leq
        \frac{1}{8n}
    \]
    が成り立つ。

    ここで，アルキメデスの性質より，
    \[
        N>\frac{1}{8\epsilon}
    \]
    を満たす自然数 $N$ を取ることができる。

    このとき，$n\geq N$ ならば
    \[
        \frac{1}{8n}
        \leq
        \frac{1}{8N}
        <
        \epsilon
    \]
    である。

    したがって，$n\geq N$ ならば
    \[
        \left|
            \sqrt{n}\left(\sqrt{n+1}-\sqrt{n}\right)
            -
            \frac{1}{2}
        \right|
        <
        \epsilon
    \]
    が成り立つ。

    ゆえに
    \[
        \lim_{n\to\infty}
        \sqrt{n}\left(\sqrt{n+1}-\sqrt{n}\right)
        =
        \frac{1}{2}
    \]
    である。
\end{proof}